% ============================================================
% ENHANCED INTERACTION STRATEGIES FOR AUTISM MATH LEARNING
% Research Report on Pedagogical Strategies and Technical Improvements
% ============================================================

\documentclass[12pt, a4paper]{article}

% ==================== PACKAGES ====================
\usepackage[utf8]{inputenc}
\usepackage[margin=1in]{geometry}
\usepackage{graphicx}
\usepackage{xcolor}
\usepackage{hyperref}
\usepackage{fancyhdr}
\usepackage{titlesec}
\usepackage{tocloft}
\usepackage{booktabs}
\usepackage{longtable}
\usepackage{array}
\usepackage{multirow}
\usepackage{enumitem}
\usepackage{amsmath}
\usepackage{amssymb}
\usepackage{float}
\usepackage{caption}
\usepackage{subcaption}
\usepackage{listings}
\usepackage{tikz}
\usepackage{tcolorbox}
\usepackage{cite}

% ==================== COLOR DEFINITIONS ====================
\definecolor{primaryblue}{RGB}{0, 102, 204}
\definecolor{secondaryblue}{RGB}{51, 153, 255}
\definecolor{accentgreen}{RGB}{0, 204, 102}
\definecolor{warningorange}{RGB}{255, 153, 51}
\definecolor{textgray}{RGB}{51, 51, 51}
\definecolor{lightgray}{RGB}{245, 245, 245}

% ==================== HYPERLINK SETUP ====================
\hypersetup{
    colorlinks=true,
    linkcolor=primaryblue,
    filecolor=magenta,      
    urlcolor=secondaryblue,
    citecolor=primaryblue,
    pdftitle={Enhanced Interaction Strategies for Autism Math Learning},
    pdfauthor={Rohith Kumar},
    pdfsubject={Autism Education Research},
    pdfkeywords={Autism, Mathematics, Interaction Design, React, Accessibility}
}

% ==================== HEADER AND FOOTER ====================
\pagestyle{fancy}
\fancyhf{}
\fancyhead[L]{\textcolor{primaryblue}{\textbf{Enhanced Interaction Strategies}}}
\fancyhead[R]{\textcolor{secondaryblue}{Research Report}}
\fancyfoot[C]{\textcolor{primaryblue}{\thepage}}
\renewcommand{\headrulewidth}{2pt}
\renewcommand{\footrulewidth}{1pt}
\renewcommand{\headrule}{\hbox to\headwidth{\color{accentgreen}\leaders\hrule height \headrulewidth\hfill}}

% ==================== TITLE FORMATTING ====================
\titleformat{\section}
{\color{primaryblue}\normalfont\Large\bfseries}
{\color{accentgreen}\thesection}{1em}{}[\color{secondaryblue}\titlerule]

\titleformat{\subsection}
{\color{secondaryblue}\normalfont\large\bfseries}
{\color{warningorange}\thesubsection}{1em}{}

\titleformat{\subsubsection}
{\color{primaryblue}\normalfont\normalsize\bfseries}
{\color{secondaryblue}\thesubsubsection}{1em}{}

% ==================== CODE LISTING STYLE ====================
\lstset{
    basicstyle=\ttfamily\small,
    breaklines=true,
    frame=single,
    backgroundcolor=\color{lightgray},
    keywordstyle=\color{primaryblue}\bfseries,
    commentstyle=\color{accentgreen},
    stringstyle=\color{warningorange},
    numbers=left,
    numberstyle=\tiny\color{textgray},
    captionpos=b
}

% ==================== DOCUMENT START ====================
\begin{document}

% ==================== TITLE PAGE ====================
\begin{titlepage}
    \centering
    \vspace*{2cm}
    
    {\Huge\bfseries\color{primaryblue} ENHANCED INTERACTION STRATEGIES\par}
    \vspace{0.5cm}
    {\LARGE\color{secondaryblue} FOR AUTISM MATH LEARNING\par}
    
    \vspace{1cm}
    
    \begin{tcolorbox}[colback=accentgreen!5, colframe=primaryblue, width=0.85\textwidth, arc=5mm]
        \centering
        {\Large\textbf{Research Report on Pedagogical Strategies}\par}
        \vspace{0.3cm}
        {\Large\textbf{and Technical Interaction Improvements}\par}
        \vspace{0.5cm}
        {\large Star Math Explorer Enhancement Project\par}
    \end{tcolorbox}
    
    \vspace{2cm}
    
    \begin{tcolorbox}[colback=secondaryblue!10, colframe=accentgreen, width=0.75\textwidth, arc=3mm]
        \centering
        \textbf{Student Information}\\
        \vspace{0.3cm}
        \textbf{Roll No:} CB.SC.U4CSE23018\\
        \textbf{Name:} Rohith Kumar\\
        \textbf{Department:} CSE-A\\
        \textbf{Institution:} Amrita Vishwa Vidyapeetham\\
        \vspace{0.3cm}
        \textbf{Course Teacher:} Dr. T. Senthil Kumar\\
        Professor, Amrita School of Computing\\
        Coimbatore - 641112
    \end{tcolorbox}
    
    \vspace{2cm}
    
    {\large GitHub Repository: \url{https://github.com/9059Rohith/lab2}\par}
    
    \vfill
    
    {\large \today\par}
\end{titlepage}

% ==================== TABLE OF CONTENTS ====================
\tableofcontents
\newpage

% ==================== ABSTRACT ====================
\section*{Abstract}
\addcontentsline{toc}{section}{Abstract}

This research report explores evidence-based pedagogical strategies for teaching mathematical concepts to children with Autism Spectrum Disorder (ASD), with specific focus on the Star Math Explorer application. The report synthesizes current research literature on autism-friendly mathematical instruction and proposes comprehensive technical enhancements to improve user interaction through advanced event handling mechanisms.

The document is divided into two major components: (1) Research-based teaching strategies specifically tailored for autism learners studying mathematical symbols and concepts, and (2) Technical implementation strategies for enhancing user interaction through keystroke events, mouse events, screen capture operations, and multimodal feedback systems using React.js ecosystem.

Key findings indicate that children with ASD benefit significantly from visual-spatial learning approaches, concrete-representational-abstract (CRA) methodology, and interactive feedback systems. The proposed technical enhancements leverage React's event handling capabilities, including keyboard navigation, gesture-based controls, screen recording for progress monitoring, and adaptive response systems that accommodate diverse sensory processing profiles.

\textbf{Keywords:} Autism Spectrum Disorder, Mathematical Learning, Interaction Design, React Event Handling, Screen Capture, Accessibility, Special Education Technology

\newpage

% ==================== SECTION 1: INTRODUCTION ====================
\section{Introduction}

\subsection{Background and Motivation}

Autism Spectrum Disorder (ASD) is a neurodevelopmental condition affecting approximately 1 in 36 children in the United States (CDC, 2023). Children with ASD often demonstrate unique cognitive profiles characterized by strengths in visual processing and pattern recognition, alongside challenges in abstract reasoning and symbolic understanding. Mathematics education presents both opportunities and obstacles for this population.

The Star Math Explorer application was developed to leverage the strengths of ASD learners while addressing their specific learning challenges through space-themed mathematical symbol instruction. However, the effectiveness of any educational technology depends not only on pedagogical soundness but also on the quality and accessibility of user interactions.

\subsection{Research Objectives}

This report addresses two interconnected objectives:

\begin{enumerate}[leftmargin=*]
    \item \textbf{Pedagogical Objective:} To identify and synthesize evidence-based strategies for teaching mathematical concepts to children with autism, with emphasis on symbol recognition, operational understanding, and conceptual development.
    
    \item \textbf{Technical Objective:} To propose and document comprehensive interaction enhancements utilizing advanced event handling mechanisms including keyboard events, mouse tracking, touch gestures, and screen capture capabilities within the React.js framework.
\end{enumerate}

\subsection{Scope and Methodology}

This research employs a literature review methodology combined with technical analysis and implementation design. Academic sources were reviewed from special education journals, autism research publications, and human-computer interaction literature. Technical specifications draw from React.js documentation, web accessibility standards (WCAG 2.1), and assistive technology best practices.

The mathematical concepts addressed in Star Math Explorer include:
\begin{itemize}[leftmargin=*]
    \item Basic operations (addition, subtraction, multiplication, division)
    \item Special mathematical symbols (pi, infinity, square root, sigma)
    \item Comparison operators (less than, greater than, equal to)
    \item Advanced concepts (squared, cubed, percentage, fractions)
\end{itemize}

\subsection{Document Organization}

This report is organized into eight sections:
\begin{itemize}[leftmargin=*]
    \item Section 2: Literature Review on Autism and Mathematical Learning
    \item Section 3: Evidence-Based Teaching Strategies for ASD Learners
    \item Section 4: Current Interaction Model Analysis
    \item Section 5: Enhanced Interaction Design Framework
    \item Section 6: React Event Handling Implementation
    \item Section 7: Screen Capture Integration for Progress Monitoring
    \item Section 8: Accessibility Compliance and Universal Design
    \item Section 9: Conclusion and Future Directions
\end{itemize}

\begin{figure}[H]
    \centering
    \includegraphics[width=0.75\textwidth]{figure-1.png}
    \caption{Portal Landing Page - Space-Themed Interface with Engaging Visual Design}
    \label{fig:landing-page}
\end{figure}

\newpage

% ==================== SECTION 2: LITERATURE REVIEW ====================
\section{Literature Review: Autism and Mathematical Learning}

\subsection{Cognitive Profile of Children with Autism}

Research consistently demonstrates that children with ASD exhibit heterogeneous cognitive profiles with specific patterns relevant to mathematical learning:

\subsubsection{Visual-Spatial Strengths}

Grandin (2009) famously described her cognitive experience as "thinking in pictures," a characteristic shared by many individuals with autism. Studies using fMRI have shown enhanced visual cortex activation in ASD populations during problem-solving tasks. This visual processing advantage can be leveraged in mathematics education through:

\begin{itemize}[leftmargin=*]
    \item Graphical representations of abstract concepts
    \item Color-coded categorical organization
    \item Spatial arrangement reflecting mathematical relationships
    \item Visual metaphors connecting symbols to concrete meanings
\end{itemize}

\subsubsection{Pattern Recognition Abilities}

Baron-Cohen et al. (2009) documented superior pattern recognition in many children with autism, particularly in structured, rule-based systems—precisely the nature of mathematics. Mathematical operations follow consistent patterns that can be visually represented, making them particularly suitable for ASD learners when properly presented.

\subsubsection{Challenges in Abstract Reasoning}

While pattern recognition may be enhanced, abstract symbolic reasoning often presents challenges. Research by Minshew et al. (2002) indicates difficulties with:

\begin{itemize}[leftmargin=*]
    \item Understanding symbols divorced from concrete referents
    \item Generalizing learning across contexts
    \item Flexible application of learned procedures
    \item Conceptual integration of related mathematical ideas
\end{itemize}

\subsection{Executive Function and Mathematical Learning}

Executive function deficits commonly associated with ASD significantly impact mathematical problem-solving:

\subsubsection{Working Memory Constraints}

Studies by Goldstein et al. (2014) reveal that many children with autism demonstrate reduced working memory capacity, particularly for verbal-sequential information. Multi-step mathematical problems requiring simultaneous information retention and manipulation pose particular challenges.

\textbf{Implication for Design:} Interfaces should minimize working memory load through:
\begin{itemize}[leftmargin=*]
    \item Persistent visual display of relevant information
    \item Step-by-step decomposition of complex operations
    \item Visual cues for procedural sequences
    \item Reduction of distracting stimuli
\end{itemize}

\subsubsection{Cognitive Flexibility Difficulties}

Mental set-shifting challenges mean that children with ASD may struggle when mathematical contexts change or when multiple solution strategies are required. However, structured learning environments with predictable patterns align well with ASD cognitive preferences.

\subsection{Sensory Processing Considerations}

Sensory processing differences in autism profoundly affect learning experiences:

\begin{itemize}[leftmargin=*]
    \item \textbf{Hypersensitivity:} Some children experience sensory overload from excessive visual stimulation, unpredictable sounds, or rapid movements
    \item \textbf{Hyposensitivity:} Others require stronger sensory input to register and attend to stimuli
    \item \textbf{Sensory Seeking:} Many engage in repetitive behaviors for sensory stimulation
\end{itemize}

\textbf{Educational Implication:} Learning platforms must provide sensory customization options, including volume control, animation speed adjustment, color theme selection, and togglable features.

\subsection{Motivation and Engagement}

Special interests and restricted focus areas characteristic of autism can be harnessed positively. Research by Winter-Messiers et al. (2007) demonstrates that incorporating special interests into curriculum significantly enhances engagement and learning outcomes. The space theme in Star Math Explorer exemplifies this approach, capitalizing on the common special interest in astronomy, space exploration, and cosmic phenomena among children with autism.

\begin{figure}[H]
    \centering
    \includegraphics[width=0.75\textwidth]{figure-2.png}
    \caption{Symbol Grid Display with Color-Coded Categories and Interactive Elements}
    \label{fig:symbol-grid}
\end{figure}

\newpage

% ==================== SECTION 3: EVIDENCE-BASED TEACHING STRATEGIES ====================
\section{Evidence-Based Teaching Strategies for Mathematical Concepts}

\subsection{Concrete-Representational-Abstract (CRA) Sequence}

The CRA instructional sequence represents best practice for teaching mathematics to students with learning differences, including autism (Witzel et al., 2003).

\subsubsection{Concrete Stage}

Learning begins with physical manipulatives and tangible objects:

\textbf{For Addition:} Physical objects (blocks, counters) are combined
\begin{itemize}[leftmargin=*]
    \item Child has 3 star tokens, receives 2 more, counts total
    \item Concrete experience creates foundational understanding
\end{itemize}

\textbf{For Multiplication:} Groups of identical quantities are assembled
\begin{itemize}[leftmargin=*]
    \item Three groups of four stars each
    \item Physical arrangement demonstrates "groups of" concept
\end{itemize}

\subsubsection{Representational Stage}

Concrete objects are replaced with visual representations:

\begin{itemize}[leftmargin=*]
    \item Pictures or icons replacing physical manipulatives
    \item Diagrams showing mathematical relationships
    \item Visual models (number lines, arrays, area models)
\end{itemize}

\textbf{Implementation in Star Math Explorer:}
The application currently emphasizes this stage through emoji symbols (🍎 + 🍎 = 🍎🍎) and visual examples. This bridges concrete and abstract effectively for ASD learners.

\subsubsection{Abstract Stage}

Mathematical symbols and notation are introduced:

\begin{itemize}[leftmargin=*]
    \item Standard symbolic notation (3 + 2 = 5)
    \item Abstract problem-solving
    \item Generalization to various contexts
\end{itemize}

\textbf{Recommendation:} The transition from representational to abstract should be gradual, with simultaneous display of both forms during the transition phase.

\subsection{Visual Supports and Graphic Organizers}

Research by Yakubova et al. (2015) demonstrates significant benefits of visual supports for students with ASD learning mathematics:

\subsubsection{Color Coding Systems}

\begin{itemize}[leftmargin=*]
    \item Consistent color associations aid memory and categorization
    \item Star Math Explorer effectively uses color (e.g., green for addition, pink for subtraction)
    \item Colors should follow intuitive associations when possible
\end{itemize}

\subsubsection{Visual Schedules and Sequences}

For multi-step problems, visual procedural supports prove invaluable:

\begin{itemize}[leftmargin=*]
    \item Step-by-step visual guides
    \item Progress indicators showing completion status
    \item Visual cues for "next action"
\end{itemize}

\subsubsection{Anchor Charts and Reference Materials}

Persistent display of key information reduces memory demands:

\begin{itemize}[leftmargin=*]
    \item Symbol reference guide always accessible
    \item Quick-reference visual aids
    \item Personalized "favorite symbols" section
\end{itemize}

\subsection{Systematic Instruction and Task Analysis}

Applied Behavior Analysis (ABA) principles inform effective autism education:

\subsubsection{Task Decomposition}

Complex skills are broken into discrete, teachable components:

\begin{enumerate}[leftmargin=*]
    \item Symbol recognition (visual discrimination)
    \item Name association (label matching)
    \item Conceptual understanding (meaning comprehension)
    \item Operational application (problem solving)
\end{enumerate}

Each component can be taught, practiced, and mastered separately before integration.

\subsubsection{Explicit Instruction}

Clear, unambiguous teaching with:
\begin{itemize}[leftmargin=*]
    \item Simple, concrete language
    \item Demonstration followed by guided practice
    \item Immediate corrective feedback
    \item Multiple exemplars for generalization
\end{itemize}

\subsubsection{Mastery-Based Progression}

Students advance only after demonstrating competence:

\begin{itemize}[leftmargin=*]
    \item Criterion-based advancement (e.g., 80\% accuracy over 3 trials)
    \item Prevents gaps in foundational knowledge
    \item Builds confidence through success experiences
\end{itemize}

\subsection{High-Intensity Practice with Varied Exemplars}

Research by Burns et al. (2010) indicates that students with autism benefit from:

\subsubsection{Distributed Practice}

\begin{itemize}[leftmargin=*]
    \item Short, frequent practice sessions rather than lengthy ones
    \item Spaced repetition schedules for retention
    \item Interleaved practice of related concepts
\end{itemize}

\subsubsection{Varied Contexts}

To promote generalization:
\begin{itemize}[leftmargin=*]
    \item Same concept presented with different representations
    \item Multiple problem formats
    \item Various real-world applications
\end{itemize}

\subsection{Technology-Enhanced Learning}

Research specifically examining computer-based mathematics instruction for autism shows promising results:

\subsubsection{Advantages of Digital Platforms}

\begin{itemize}[leftmargin=*]
    \item \textbf{Consistency:} Predictable, reliable responses reduce anxiety
    \item \textbf{Patience:} Unlimited practice without frustration
    \item \textbf{Immediate Feedback:} Instant reinforcement or correction
    \item \textbf{Customization:} Adaptive pacing and sensory preferences
    \item \textbf{Data Collection:} Automatic progress monitoring
\end{itemize}

\subsubsection{Interactive Multimedia}

Studies by Pennington (2010) demonstrate enhanced engagement through:
\begin{itemize}[leftmargin=*]
    \item Animation showing dynamic mathematical relationships
    \item Audio support for reading comprehension
    \item Gamification elements for motivation
    \item Interactive manipulation of visual elements
\end{itemize}

\subsection{Social Stories and Mathematical Narratives}

Gray's Social Stories methodology can be adapted for mathematical concepts:

\begin{itemize}[leftmargin=*]
    \item Story contexts making abstract symbols meaningful
    \item Character-driven scenarios applying mathematics
    \item Predictable narratives incorporating computational practice
\end{itemize}

\textbf{Example:} "Stella the Space Explorer needs to combine 3 red stars and 2 blue stars. She uses the PLUS symbol (+) to add them together. Now she has 5 stars total!"

\subsection{Multisensory Integration}

Engaging multiple sensory modalities enhances learning:

\begin{itemize}[leftmargin=*]
    \item \textbf{Visual:} Symbols, colors, animations, spatial layouts
    \item \textbf{Auditory:} Pronunciation, explanations, musical cues
    \item \textbf{Kinesthetic:} Interactive clicking, dragging, drawing
    \item \textbf{Haptic:} Vibration feedback (on touch devices)
\end{itemize}

Research by Schlosser et al. (2012) indicates particular benefit from paired visual-auditory presentation for symbol learning in autism populations.

\newpage

% ==================== SECTION 4: CURRENT INTERACTION MODEL ANALYSIS ====================
\section{Current Interaction Model Analysis}

\subsection{Existing Interaction Mechanisms}

The current Star Math Explorer implementation includes several interaction paradigms:

\subsubsection{Click-Based Interactions}

\begin{lstlisting}[language=JavaScript, caption=Current Click Handler]
const handleSymbolClick = (symbol) => {
  if (voiceEnabled) {
    speakSymbol(symbol.name);
  }
  // Trigger visual feedback
  setSelectedSymbol(symbol);
};
\end{lstlisting}

\textbf{Strengths:}
\begin{itemize}[leftmargin=*]
    \item Simple, intuitive interaction model
    \item Accessible across device types
    \item Clear cause-effect relationship
\end{itemize}

\textbf{Limitations:}
\begin{itemize}[leftmargin=*]
    \item Single interaction modality
    \item No keyboard accessibility for motor-impaired users
    \item Limited engagement variety
\end{itemize}

\subsubsection{Voice Toggle System}

Users can enable/disable audio feedback globally.

\textbf{Strengths:}
\begin{itemize}[leftmargin=*]
    \item Sensory customization
    \item User control over auditory input
\end{itemize}

\textbf{Limitations:}
\begin{itemize}[leftmargin=*]
    \item Binary on/off lacks granular control
    \item No volume adjustment
    \item No per-symbol audio preferences
\end{itemize}

\subsubsection{Visual Feedback Systems}

Confetti animations and color changes provide positive reinforcement.

\textbf{Strengths:}
\begin{itemize}[leftmargin=*]
    \item Engaging visual rewards
    \item Immediate feedback
\end{itemize}

\textbf{Limitations:}
\begin{itemize}[leftmargin=*]
    \item May be overstimulating for some users
    \item No customization of animation intensity
    \item Limited variety in feedback types
\end{itemize}

\subsection{Identified Gaps and Enhancement Opportunities}

\begin{enumerate}[leftmargin=*]
    \item \textbf{Keyboard Navigation:} No keyboard-only operation support
    \item \textbf{Mouse Event Diversity:} Only click events, no hover, drag, or gesture support
    \item \textbf{Progress Documentation:} No screen capture or recording functionality
    \item \textbf{Adaptive Feedback:} No personalization based on user performance
    \item \textbf{Alternative Input Methods:} No voice input or gesture recognition
    \item \textbf{Attention Monitoring:} No detection of engagement levels
    \item \textbf{Error Prevention:} Limited guidance to prevent incorrect interactions
\end{enumerate}

\newpage

% ==================== SECTION 5: ENHANCED INTERACTION DESIGN FRAMEWORK ====================
\section{Enhanced Interaction Design Framework}

\subsection{Multimodal Interaction Architecture}

The enhanced system implements multiple, complementary interaction channels:

\begin{tcolorbox}[colback=lightgray, colframe=primaryblue, title=Interaction Modalities]
\begin{itemize}[leftmargin=*]
    \item \textbf{Keyboard Events:} Arrow navigation, Enter selection, number shortcuts
    \item \textbf{Mouse Events:} Click, double-click, hover, right-click context menus
    \item \textbf{Touch/Gesture:} Tap, swipe, pinch-zoom, long-press
    \item \textbf{Screen Events:} Scroll tracking, visibility detection, window focus
    \item \textbf{Voice Input:} Speech recognition for symbol selection
    \item \textbf{Screen Capture:} Recording sessions for review and assessment
\end{itemize}
\end{tcolorbox}

\subsection{Event Handling Hierarchy}

React's synthetic event system provides cross-browser consistency:

\begin{lstlisting}[language=JavaScript, caption=Event Handler Architecture]
// Hierarchical event management
const eventHandlers = {
  keyboard: useKeyboardNavigation(),
  mouse: useMouseInteractions(),
  touch: useTouchGestures(),
  screen: useScreenMonitoring(),
  capture: useScreenCapture()
};
\end{lstlisting}

\subsection{Accessibility-First Design Principles}

Following WCAG 2.1 Level AA standards:

\begin{itemize}[leftmargin=*]
    \item \textbf{Perceivable:} Multiple sensory channels for information
    \item \textbf{Operable:} Multiple navigation methods available
    \item \textbf{Understandable:} Predictable, consistent interactions
    \item \textbf{Robust:} Compatible with assistive technologies
\end{itemize}

\subsection{Adaptive Response System}

Machine learning-informed adaptation:

\begin{enumerate}[leftmargin=*]
    \item Monitor user interaction patterns
    \item Identify preferences and difficulties
    \item Adjust interface responsiveness
    \item Personalize feedback timing and intensity
\end{enumerate}

\newpage

% ==================== SECTION 6: REACT EVENT HANDLING IMPLEMENTATION ====================
\section{React Event Handling Implementation}

\subsection{Keyboard Event Integration}

\subsubsection{Arrow Key Navigation}

Enables keyboard-only symbol browsing crucial for accessibility:

\begin{lstlisting}[language=JavaScript, caption=Keyboard Navigation Hook]
import { useEffect, useState } from 'react';

const useKeyboardNavigation = (symbols) => {
  const [focusedIndex, setFocusedIndex] = useState(0);
  
  useEffect(() => {
    const handleKeyDown = (event) => {
      switch(event.key) {
        case 'ArrowRight':
          event.preventDefault();
          setFocusedIndex(prev => 
            (prev + 1) % symbols.length
          );
          break;
        case 'ArrowLeft':
          event.preventDefault();
          setFocusedIndex(prev => 
            (prev - 1 + symbols.length) % symbols.length
          );
          break;
        case 'ArrowDown':
          event.preventDefault();
          setFocusedIndex(prev => 
            Math.min(prev + 4, symbols.length - 1)
          );
          break;
        case 'ArrowUp':
          event.preventDefault();
          setFocusedIndex(prev => 
            Math.max(prev - 4, 0)
          );
          break;
        case 'Enter':
        case ' ':
          event.preventDefault();
          handleSymbolSelect(symbols[focusedIndex]);
          break;
        default:
          break;
      }
    };
    
    window.addEventListener('keydown', handleKeyDown);
    return () => window.removeEventListener('keydown', handleKeyDown);
  }, [symbols, focusedIndex]);
  
  return focusedIndex;
};
\end{lstlisting}

\subsubsection{Number Shortcut Keys}

Direct access to symbols via numeric keys:

\begin{lstlisting}[language=JavaScript, caption=Numeric Shortcuts]
const handleNumberShortcut = (event) => {
  const num = parseInt(event.key);
  if (!isNaN(num) && num >= 1 && num <= 9) {
    const symbolIndex = num - 1;
    if (symbolIndex < visibleSymbols.length) {
      handleSymbolSelect(visibleSymbols[symbolIndex]);
      speakSymbol(visibleSymbols[symbolIndex].name);
    }
  }
};
\end{lstlisting}

\subsubsection{Escape Key for Exit/Reset}

\begin{lstlisting}[language=JavaScript, caption=Escape Key Handler]
useEffect(() => {
  const handleEscape = (event) => {
    if (event.key === 'Escape') {
      setSelectedSymbol(null);
      setShowDetails(false);
      // Return focus to main navigation
      document.getElementById('main-nav')?.focus();
    }
  };
  
  document.addEventListener('keydown', handleEscape);
  return () => document.removeEventListener('keydown', handleEscape);
}, []);
\end{lstlisting}

\begin{figure}[H]
    \centering
    \includegraphics[width=0.7\textwidth]{figure-3.png}
    \caption{Keyboard Navigation System with Visual Focus Indicators}
    \label{fig:keyboard-nav}
\end{figure}

\subsection{Enhanced Mouse Event Handling}

\subsubsection{Hover Effects with Delayed Information Display}

\begin{lstlisting}[language=JavaScript, caption=Smart Hover System]
const useSmartHover = (symbol) => {
  const [isHovered, setIsHovered] = useState(false);
  const [showTooltip, setShowTooltip] = useState(false);
  const hoverTimerRef = useRef(null);
  
  const handleMouseEnter = () => {
    setIsHovered(true);
    // Delay tooltip to avoid overwhelming quick scans
    hoverTimerRef.current = setTimeout(() => {
      setShowTooltip(true);
      if (voiceEnabled) {
        speakSymbol(symbol.name, { rate: 0.9 });
      }
    }, 800); // 800ms delay prevents accidental triggers
  };
  
  const handleMouseLeave = () => {
    setIsHovered(false);
    setShowTooltip(false);
    if (hoverTimerRef.current) {
      clearTimeout(hoverTimerRef.current);
    }
  };
  
  return { isHovered, showTooltip, handleMouseEnter, handleMouseLeave };
};
\end{lstlisting}

\subsubsection{Right-Click Context Menu}

\begin{lstlisting}[language=JavaScript, caption=Context Menu Implementation]
const handleContextMenu = (event, symbol) => {
  event.preventDefault();
  
  setContextMenu({
    visible: true,
    x: event.clientX,
    y: event.clientY,
    symbol: symbol,
    options: [
      { label: 'Hear pronunciation', action: () => speakSymbol(symbol.name) },
      { label: 'See examples', action: () => showExamples(symbol) },
      { label: 'Practice problems', action: () => startPractice(symbol) },
      { label: 'Add to favorites', action: () => addToFavorites(symbol) }
    ]
  });
};
\end{lstlisting}

\subsubsection{Double-Click for Detailed View}

\begin{lstlisting}[language=JavaScript, caption=Double-Click Handler]
const handleDoubleClick = (symbol) => {
  setDetailedViewSymbol(symbol);
  setShowDetailedModal(true);
  
  // Speak detailed description
  if (voiceEnabled) {
    speakSymbol(`${symbol.name}. ${symbol.description}`);
  }
  
  // Track engagement analytics
  logInteraction('double_click_detail_view', symbol.symbol);
};
\end{lstlisting}

\subsubsection{Mouse Movement Tracking for Engagement}

\begin{lstlisting}[language=JavaScript, caption=Engagement Detection]
const useEngagementTracking = () => {
  const [lastActivity, setLastActivity] = useState(Date.now());
  const [isEngaged, setIsEngaged] = useState(true);
  
  useEffect(() => {
    const handleActivity = () => {
      setLastActivity(Date.now());
      setIsEngaged(true);
    };
    
    const checkEngagement = setInterval(() => {
      const idleTime = Date.now() - lastActivity;
      if (idleTime > 30000) { // 30 seconds idle
        setIsEngaged(false);
        // Trigger gentle re-engagement prompt
        showReengagementPrompt();
      }
    }, 5000);
    
    window.addEventListener('mousemove', handleActivity);
    window.addEventListener('click', handleActivity);
    window.addEventListener('keydown', handleActivity);
    
    return () => {
      clearInterval(checkEngagement);
      window.removeEventListener('mousemove', handleActivity);
      window.removeEventListener('click', handleActivity);
      window.removeEventListener('keydown', handleActivity);
    };
  }, [lastActivity]);
  
  return isEngaged;
};
\end{lstlisting}

\subsection{Touch and Gesture Events}

\subsubsection{Swipe Gestures for Navigation}

\begin{lstlisting}[language=JavaScript, caption=Swipe Detection]
const useSwipeGesture = (onSwipeLeft, onSwipeRight) => {
  const [touchStart, setTouchStart] = useState(null);
  const [touchEnd, setTouchEnd] = useState(null);
  
  const minSwipeDistance = 50;
  
  const onTouchStart = (e) => {
    setTouchEnd(null);
    setTouchStart(e.targetTouches[0].clientX);
  };
  
  const onTouchMove = (e) => {
    setTouchEnd(e.targetTouches[0].clientX);
  };
  
  const onTouchEnd = () => {
    if (!touchStart || !touchEnd) return;
    
    const distance = touchStart - touchEnd;
    const isLeftSwipe = distance > minSwipeDistance;
    const isRightSwipe = distance < -minSwipeDistance;
    
    if (isLeftSwipe && onSwipeLeft) {
      onSwipeLeft();
      provideTactileFeedback();
    }
    if (isRightSwipe && onSwipeRight) {
      onSwipeRight();
      provideTactileFeedback();
    }
  };
  
  return { onTouchStart, onTouchMove, onTouchEnd };
};

const provideTactileFeedback = () => {
  if ('vibrate' in navigator) {
    navigator.vibrate(50); // 50ms vibration
  }
};
\end{lstlisting}

\subsubsection{Long-Press for Additional Options}

\begin{lstlisting}[language=JavaScript, caption=Long-Press Detection]
const useLongPress = (callback, ms = 500) => {
  const [startLongPress, setStartLongPress] = useState(false);
  
  useEffect(() => {
    let timerId;
    if (startLongPress) {
      timerId = setTimeout(() => {
        callback();
        provideTactileFeedback();
      }, ms);
    } else {
      clearTimeout(timerId);
    }
    
    return () => {
      clearTimeout(timerId);
    };
  }, [startLongPress, callback, ms]);
  
  return {
    onMouseDown: () => setStartLongPress(true),
    onMouseUp: () => setStartLongPress(false),
    onMouseLeave: () => setStartLongPress(false),
    onTouchStart: () => setStartLongPress(true),
    onTouchEnd: () => setStartLongPress(false),
  };
};
\end{lstlisting}

\begin{figure}[H]
    \centering
    \includegraphics[width=0.7\textwidth]{figure-4.png}
    \caption{Touch Gesture Interactions and Multi-Sensory Feedback System}
    \label{fig:touch-gestures}
\end{figure}

\subsection{Screen Event Monitoring}

\subsubsection{Intersection Observer for Visibility Tracking}

\begin{lstlisting}[language=JavaScript, caption=Visibility Detection]
const useVisibilityTracking = (elementRef, onVisible) => {
  useEffect(() => {
    const observer = new IntersectionObserver(
      ([entry]) => {
        if (entry.isIntersecting) {
          onVisible(entry.target);
        }
      },
      {
        threshold: 0.5, // 50% visible
        rootMargin: '0px'
      }
    );
    
    if (elementRef.current) {
      observer.observe(elementRef.current);
    }
    
    return () => {
      if (elementRef.current) {
        observer.unobserve(elementRef.current);
      }
    };
  }, [elementRef, onVisible]);
};

// Usage for auto-announcement
const SymbolCard = ({ symbol }) => {
  const cardRef = useRef(null);
  
  useVisibilityTracking(cardRef, () => {
    if (voiceEnabled && autoAnnounce) {
      speakSymbol(symbol.name);
    }
  });
  
  return <div ref={cardRef} className="symbol-card">...</div>;
};
\end{lstlisting}

\subsubsection{Scroll Event Handling}

\begin{lstlisting}[language=JavaScript, caption=Scroll Progress Tracking]
const useScrollProgress = () => {
  const [scrollProgress, setScrollProgress] = useState(0);
  
  useEffect(() => {
    const handleScroll = () => {
      const windowHeight = window.innerHeight;
      const documentHeight = document.documentElement.scrollHeight;
      const scrollTop = window.pageYOffset || document.documentElement.scrollTop;
      
      const progress = (scrollTop / (documentHeight - windowHeight)) * 100;
      setScrollProgress(Math.min(100, Math.max(0, progress)));
    };
    
    window.addEventListener('scroll', handleScroll, { passive: true });
    return () => window.removeEventListener('scroll', handleScroll);
  }, []);
  
  return scrollProgress;
};
\end{lstlisting}

\subsubsection{Window Focus Detection}

\begin{lstlisting}[language=JavaScript, caption=Focus Management]
const useWindowFocus = () => {
  const [isFocused, setIsFocused] = useState(true);
  
  useEffect(() => {
    const handleFocus = () => {
      setIsFocused(true);
      // Resume animations and audio
    };
    
    const handleBlur = () => {
      setIsFocused(false);
      // Pause background animations to save resources
      // Pause audio to prevent disruption
    };
    
    window.addEventListener('focus', handleFocus);
    window.addEventListener('blur', handleBlur);
    
    return () => {
      window.removeEventListener('focus', handleFocus);
      window.removeEventListener('blur', handleBlur);
    };
  }, []);
  
  return isFocused;
};
\end{lstlisting}

\newpage

% ==================== SECTION 7: SCREEN CAPTURE INTEGRATION ====================
\section{Screen Capture Integration for Progress Monitoring}

\subsection{Rationale for Screen Recording in Special Education}

Screen capture functionality serves multiple pedagogical purposes in autism education:

\begin{itemize}[leftmargin=*]
    \item \textbf{Progress Documentation:} Visual evidence of skill development
    \item \textbf{Behavioral Analysis:} Review interaction patterns and challenges
    \item \textbf{Parent Communication:} Share learning sessions with family
    \item \textbf{IEP Documentation:} Data for Individualized Education Programs
    \item \textbf{Self-Review:} Children can watch their own learning journey
    \item \textbf{Instructional Refinement:} Identify areas needing modification
\end{itemize}

\subsection{React Screen Capture Implementation}

\subsubsection{Using react-screen-capture Library}

\begin{lstlisting}[language=JavaScript, caption=Screen Capture Setup]
import { useScreenCapture } from 'react-screen-capture';

const ScreenRecordingComponent = () => {
  const {
    startRecording,
    stopRecording,
    mediaBlobUrl,
    status,
    error
  } = useScreenCapture();
  
  return (
    <div className="recording-controls">
      {status === 'idle' && (
        <button onClick={startRecording}>
          Start Recording Session
        </button>
      )}
      {status === 'recording' && (
        <button onClick={stopRecording}>
          Stop Recording
        </button>
      )}
      {mediaBlobUrl && (
        <div>
          <video src={mediaBlobUrl} controls />
          <button onClick={() => downloadRecording(mediaBlobUrl)}>
            Download Session
          </button>
        </div>
      )}
    </div>
  );
};
\end{lstlisting}

\subsubsection{Custom Screen Recording Hook}

\begin{lstlisting}[language=JavaScript, caption=Custom Recording Hook]
import { useState, useRef, useCallback } from 'react';

const useCustomScreenRecording = () => {
  const [recording, setRecording] = useState(false);
  const [recordedChunks, setRecordedChunks] = useState([]);
  const mediaRecorderRef = useRef(null);
  const streamRef = useRef(null);
  
  const startRecording = useCallback(async () => {
    try {
      const stream = await navigator.mediaDevices.getDisplayMedia({
        video: {
          width: 1920,
          height: 1080,
          frameRate: 30
        },
        audio: {
          echoCancellation: true,
          noiseSuppression: true,
          sampleRate: 44100
        }
      });
      
      streamRef.current = stream;
      
      const mediaRecorder = new MediaRecorder(stream, {
        mimeType: 'video/webm;codecs=vp9',
        videoBitsPerSecond: 2500000
      });
      
      mediaRecorder.ondataavailable = (event) => {
        if (event.data.size > 0) {
          setRecordedChunks(prev => [...prev, event.data]);
        }
      };
      
      mediaRecorder.onstop = () => {
        stream.getTracks().forEach(track => track.stop());
      };
      
      mediaRecorderRef.current = mediaRecorder;
      mediaRecorder.start(100); // Collect data every 100ms
      setRecording(true);
      
      // Log recording start
      console.log('Recording started:', new Date().toISOString());
      
    } catch (error) {
      console.error('Error starting recording:', error);
      alert('Unable to start recording. Please check permissions.');
    }
  }, []);
  
  const stopRecording = useCallback(() => {
    if (mediaRecorderRef.current && recording) {
      mediaRecorderRef.current.stop();
      setRecording(false);
      
      // Log recording stop
      console.log('Recording stopped:', new Date().toISOString());
    }
  }, [recording]);
  
  const downloadRecording = useCallback(() => {
    if (recordedChunks.length === 0) return;
    
    const blob = new Blob(recordedChunks, { type: 'video/webm' });
    const url = URL.createObjectURL(blob);
    const a = document.createElement('a');
    const timestamp = new Date().toI() + '.webm';
    a.style.display = 'none';
    a.href = url;
    a.download = filename;
    document.body.appendChild(a);
    a.click();
    
    // Cleanup
    setTimeout(() => {
      document.body.removeChild(a);
      URL.revokeObjectURL(url);
    }, 100);
    
    // Clear recorded chunks
    setRecordedChunks([]);
  }, [recordedChunks]);
  
  const uploadRecording = useCallback(async (childId) => {
    if (recordedChunks.length === 0) return;
    
    const blob = new Blob(recordedChunks, { type: 'video/webm' });
    const formData = new FormData();
    formData.append('video', blob, `session_${Date.now()}.webm`);
    formData.append('childId', childId);
    formData.append('timestamp', new Date().toISOString());
    
    try {
      const response = await fetch('/api/upload-session', {
        method: 'POST',
        body: formData
      });
      
      if (response.ok) {
        console.log('Session uploaded successfully');
        setRecordedChunks([]);
      }
    } catch (error) {
      console.error('Upload failed:', error);
    }
  }, [recordedChunks]);
  
  return {
    recording,
    startRecording,
    stopRecording,
    downloadRecording,
    uploadRecording,
    hasRecording: recordedChunks.length > 0
  };
};
\end{lstlisting}

\subsection{Privacy and Consent Considerations}

\begin{tcolorbox}[colback=warningorange!10, colframe=warningorange, title=Important Privacy Guidelines]
\textbf{Legal and Ethical Requirements:}
\begin{itemize}[leftmargin=*]
    \item Obtain explicit parental consent before recording
    \item Display clear recording indicator at all times
    \item Store recordings securely with encryption
    \item Provide easy deletion mechanism
    \item Comply with FERPA (Family Educational Rights and Privacy Act)
    \item Implement COPPA (Children's Online Privacy Protection Act) compliance
    \item Allow parents to control recording features
\end{itemize}
\end{tcolorbox}

\subsection{Automated Highlights and Key Moments}

\begin{lstlisting}[language=JavaScript, caption=Auto-highlight Detection]
const useKeyMomentDetection = () => {
  const [keyMoments, setKeyMoments] = useState([]);
  
  const recordKeyMoment = useCallback((type, data) => {
    const moment = {
      timestamp: Date.now(),
      type: type, // 'first_success', 'mastery', 'struggle', 'breakthrough'
      data: data,
      videoTimestamp: getCurrentRecordingTime()
    };
    
    setKeyMoments(prev => [...prev, moment]);
  }, []);
  
  // Detect significant events
  useEffect(() => {
    // First time correctly identifying a symbol
    if (correctAnswers === 1) {
      recordKeyMoment('first_success', { symbol: currentSymbol });
    }
    
    // Achieving mastery (e.g., 5 consecutive correct)
    if (consecutiveCorrect === 5) {
      recordKeyMoment('mastery', { symbol: currentSymbol });
    }
    
    // Struggle detection (multiple incorrect attempts)
    if (incorrectAttempts >= 3) {
      recordKeyMoment('struggle', { symbol: currentSymbol });
    }
  }, [correctAnswers, consecutiveCorrect, incorrectAttempts]);
  
  return { keyMoments, recordKeyMoment };
};
\end{lstlisting}

\subsection{Integration with Learning Management}

\begin{lstlisting}[language=JavaScript, caption=Session Data Structure]
const LearningSession = {
  sessionId: 'uuid-v4',
  childId: 'student-identifier',
  startTime: '2026-02-19T10:30:00Z',
  endTime: '2026-02-19T10:45:00Z',
  duration: 900, // seconds
  symbolsViewed: ['plus', 'minus', 'multiply'],
  symbolsLearned: ['plus', 'minus'],
  interactions: [
    { timestamp: '10:31:05', type: 'click', symbol: 'plus' },
    { timestamp: '10:31:12', type: 'hover', symbol: 'minus' },
    // ... more interactions
  ],
  keyMoments: [
    { timestamp: '10:32:00', type: 'first_success', symbol: 'plus' }
  ],
  videoUrl: 'blob:https://example.com/session-recording',
  performanceMetrics: {
    attentionScore: 85,
    engagementLevel: 'high',
    frustractionEvents: 1,
    successRate: 0.75
  }
};
\end{lstlisting}

\begin{figure}[H]
    \centering
    \includegraphics[width=0.75\textwidth]{figure-5.png}
    \caption{Screen Recording Interface with Progress Tracking and Key Moments Highlighting}
    \label{fig:screen-recording}
\end{figure}

\newpage

% ==================== SECTION 8: ACCESSIBILITY AND UNIVERSAL DESIGN ====================
\section{Accessibility Compliance and Universal Design}

\subsection{Web Content Accessibility Guidelines (WCAG) 2.1}

\subsubsection{Perceivable Information}

\textbf{1.1 Text Alternatives:}
\begin{itemize}[leftmargin=*]
    \item All images have descriptive alt text
    \item Mathematical symbols include aria-label descriptions
    \item Decorative elements marked with empty alt=""
\end{itemize}

\begin{lstlisting}[language=JavaScript, caption=Accessible Symbol Component]
const AccessibleSymbol = ({ symbol }) => (
  <div
    role="button"
    tabIndex={0}
    aria-label={`${symbol.name}. ${symbol.description}`}
    aria-describedby={`desc-${symbol.symbol}`}
    onClick={() => handleSelect(symbol)}
    onKeyPress={(e) => e.key === 'Enter' && handleSelect(symbol)}
  >
    <span className="symbol" aria-hidden="true">
      {symbol.symbol}
    </span>
    <span id={`desc-${symbol.symbol}`} className="sr-only">
      {symbol.description}
    </span>
  </div>
);
\end{lstlisting}

\textbf{1.3 Adaptable Content:}
\begin{itemize}[leftmargin=*]
    \item Semantic HTML structure
    \item Proper heading hierarchy
    \item ARIA landmarks for navigation
\end{itemize}

\textbf{1.4 Distinguishable:}
\begin{itemize}[leftmargin=*]
    \item Color contrast ratio ≥ 4.5:1 for normal text
    \item Color contrast ratio ≥ 3:1 for large text and UI components
    \item Information not conveyed by color alone
    \item Text resizable up to 200\% without loss of functionality
\end{itemize}

\subsubsection{Operable Interface}

\textbf{2.1 Keyboard Accessible:}
\begin{itemize}[leftmargin=*]
    \item All functionality available via keyboard
    \item No keyboard traps
    \item Visible focus indicators
    \item Logical tab order
\end{itemize}

\textbf{2.2 Enough Time:}
\begin{itemize}[leftmargin=*]
    \item No time limits on interactions
    \item Animations can be paused
    \item Auto-playing audio can be stopped
\end{itemize}

\textbf{2.3 Seizures and Physical Reactions:}
\begin{itemize}[leftmargin=*]
    \item No content flashing more than 3 times per second
    \item Animation motion can be disabled
\end{itemize}

\textbf{2.4 Navigable:}
\begin{itemize}[leftmargin=*]
    \item Skip navigation links
    \item Descriptive page titles
    \item Focus order follows reading order
    \item Link purpose clear from context
\end{itemize}

\textbf{2.5 Input Modalities:}
\begin{itemize}[leftmargin=*]
    \item Gestures have single-pointer alternatives
    \item Touch targets minimum 44x44 pixels
    \item Label text included in accessible name
\end{itemize}

\subsubsection{Understandable Content}

\textbf{3.1 Readable:}
\begin{itemize}[leftmargin=*]
    \item Language specified (lang="en")
    \item Simple, clear language appropriate for cognitive level
    \item Definitions for unusual terms
\end{itemize}

\textbf{3.2 Predictable:}
\begin{itemize}[leftmargin=*]
    \item Consistent navigation patterns
    \item Consistent identification of components
    \item Changes in context only when expected
\end{itemize}

\textbf{3.3 Input Assistance:}
\begin{itemize}[leftmargin=*]
    \item Error identification and description
    \item Labels or instructions for inputs
    \item Error prevention mechanisms
\end{itemize}

\subsection{Autism-Specific Accessibility Enhancements}

\subsubsection{Sensory Customization Panel}

\begin{lstlisting}[language=JavaScript, caption=Sensory Preferences System]
const SensoryControlPanel = () => {
  const [preferences, setPreferences] = useState({
    animations: true,
    animationSpeed: 1.0,
    soundEffects: true,
    volume: 0.7,
    speechRate: 1.0,
    backgroundMovement: true,
    colorScheme: 'default', // 'default', 'high-contrast', 'muted'
    particleDensity: 'medium' // 'low', 'medium', 'high', 'off'
  });
  
  return (
    <div className="sensory-controls">
      <h2>Customize Your Experience</h2>
      
      <label>
        Animation Speed:
        <input
          type="range"
          min="0.5"
          max="2"
          step="0.1"
          value={preferences.animationSpeed}
          onChange={(e) => updatePreference('animationSpeed', e.target.value)}
        />
      </label>
      
      <label>
        Sound Volume:
        <input
          type="range"
          min="0"
          max="1"
          step="0.1"
          value={preferences.volume}
          onChange={(e) => updatePreference('volume', e.target.value)}
        />
      </label>
      
      <label>
        Color Scheme:
        <select
          value={preferences.colorScheme}
          onChange={(e) => updatePreference('colorScheme', e.target.value)}
        >
          <option value="default">Colorful (Default)</option>
          <option value="high-contrast">High Contrast</option>
          <option value="muted">Calm Colors</option>
        </select>
      </label>
      
      {/* More controls... */}
    </div>
  );
};
\end{lstlisting}

\subsubsection{Reduce Motion Support}

\begin{lstlisting}[language=JavaScript, caption=Respecting prefers-reduced-motion]
const useReducedMotion = () => {
  const [prefersReducedMotion, setPrefersReducedMotion] = useState(false);
  
  useEffect(() => {
    const mediaQuery = window.matchMedia('(prefers-reduced-motion: reduce)');
    setPrefersReducedMotion(mediaQuery.matches);
    
    const handleChange = () => {
      setPrefersReducedMotion(mediaQuery.matches);
    };
    
    mediaQuery.addEventListener('change', handleChange);
    return () => mediaQuery.removeEventListener('change', handleChange);
  }, []);
  
  return prefersReducedMotion;
};

// Usage in component
const AnimatedElement = ({ children }) => {
  const prefersReducedMotion = useReducedMotion();
  
  const animationClass = prefersReducedMotion
    ? 'no-animation'
    : 'with-animation';
  
  return (
    <div className={animationClass}>
      {children}
    </div>
  );
};
\end{lstlisting}

\subsection{Assistive Technology Compatibility}

\subsubsection{Screen Reader Optimization}

\begin{itemize}[leftmargin=*]
    \item ARIA live regions announce dynamic content
    \item Status messages announced without moving focus
    \item Progress updates communicated clearly
    \item Form validation errors announced immediately
\end{itemize}

\begin{lstlisting}[language=JavaScript, caption=Live Region Announcements]
const LiveAnnouncer = () => {
  const [announcement, setAnnouncement] = useState('');
  
  const announce = (message, politeness = 'polite') => {
    setAnnouncement(''); // Clear first
    setTimeout(() => {
      setAnnouncement(message);
    }, 100);
  };
  
  return (
    <div
      role="status"
      aria-live="polite"
      aria-atomic="true"
      className="sr-only"
    >
      {announcement}
    </div>
  );
};
\end{lstlisting}

\subsubsection{Switch Control Support}

For children with motor difficulties using switch access:

\begin{lstlisting}[language=JavaScript, caption=Switch Scanning Implementation]
const useSwitchScanning = (items) => {
  const [scanIndex, setScanIndex] = useState(0);
  const [isScanning, setIsScanning] = useState(false);
  const scanIntervalRef = useRef(null);
  
  const startScanning = () => {
    setIsScanning(true);
    scanIntervalRef.current = setInterval(() => {
      setScanIndex(prev => (prev + 1) % items.length);
      playHighlightSound();
    }, 1500); // 1.5 seconds per item
  };
  
  const stopScanning = () => {
    if (scanIntervalRef.current) {
      clearInterval(scanIntervalRef.current);
    }
    setIsScanning(false);
  };
  
  const selectCurrentItem = () => {
    handleItemSelection(items[scanIndex]);
    provideFeedback('selected', items[scanIndex]);
  };
  
  return { scanIndex, startScanning, stopScanning, selectCurrentItem };
};
\end{lstlisting}

\newpage

% ==================== SECTION 9: CONCLUSION ====================
\section{Conclusion and Future Directions}

\subsection{Summary of Key Findings}

This research report has synthesized evidence-based pedagogical strategies for teaching mathematical concepts to children with Autism Spectrum Disorder and proposed comprehensive technical enhancements to the Star Math Explorer application. The key findings include:

\subsubsection{Pedagogical Insights}

\begin{enumerate}[leftmargin=*]
    \item \textbf{Visual-Spatial Learning:} Children with ASD demonstrate superior visual processing abilities that can be leveraged through graphical mathematical representations, color coding, and spatial organization.
    
    \item \textbf{CRA Methodology:} The Concrete-Representational-Abstract sequence provides optimal scaffolding for mathematical concept development, with particular emphasis on extended time in representational stages.
    
    \item \textbf{Systematic Instruction:} Task analysis, explicit teaching, and mastery-based progression align with ASD cognitive profiles and learning preferences.
    
    \item \textbf{Multisensory Integration:} Combining visual, auditory, and kinesthetic modalities strengthens memory formation and engagement while accommodating diverse sensory processing profiles.
    
    \item \textbf{Technology-Enhanced Learning:} Digital platforms offer consistency, patience, immediate feedback, and adaptability that particularly benefit autism learners.
\end{enumerate}

\subsubsection{Technical Enhancements}

\begin{enumerate}[leftmargin=*]
    \item \textbf{Multimodal Interaction:} Implementation of keyboard, mouse, touch, and gesture events provides multiple access pathways accommodating diverse motor abilities and preferences.
    
    \item \textbf{Event-Driven Engagement:} Advanced event handling enables attention monitoring, adaptive feedback, and personalized response systems.
    
    \item \textbf{Screen Capture Integration:} Recording capabilities support progress documentation, behavioral analysis, and parental involvement in learning processes.
    
    \item \textbf{Accessibility Compliance:} WCAG 2.1 adherence ensures usability for diverse abilities including screen reader users, keyboard-only navigation, and switch control access.
    
    \item \textbf{Sensory Customization:} Granular control over animations, sounds, colors, and motion respects individual sensory processing differences.
\end{enumerate}

\subsection{Implementation Priorities}

Based on impact and feasibility analysis, the following implementation priority is recommended:

\begin{tcolorbox}[colback=accentgreen!5, colframe=primaryblue, title=Phase 1: Core Accessibility (Weeks 1-2)]
\begin{itemize}[leftmargin=*]
    \item Keyboard navigation system
    \item ARIA labels and semantic structure
    \item Focus management
    \item Screen reader optimization
\end{itemize}
\end{tcolorbox}

\begin{tcolorbox}[colback=secondaryblue!5, colframe=accentgreen, title=Phase 2: Enhanced Interactions (Weeks 3-4)]
\begin{itemize}[leftmargin=*]
    \item Mouse hover effects with delayed information
    \item Touch gesture support (swipe, long-press)
    \item Engagement tracking
    \item Context menus for advanced options
\end{itemize}
\end{tcolorbox}

\begin{tcolorbox}[colback=warningorange!5, colframe=secondaryblue, title=Phase 3: Recording \& Analytics (Weeks 5-6)]
\begin{itemize}[leftmargin=*]
    \item Screen capture integration using react-screen-capture
    \item Key moment detection
    \item Session data structure and storage
    \item Privacy consent system
\end{itemize}
\end{tcolorbox}

\begin{tcolorbox}[colback=lightgray, colframe=primaryblue, title=Phase 4: Personalization (Weeks 7-8)]
\begin{itemize}[leftmargin=*]
    \item Sensory customization panel
    \item User preference persistence
    \item Adaptive difficulty adjustment
    \item Learning path optimization
\end{itemize}
\end{tcolorbox}

\subsection{Expected Outcomes}

Implementation of these enhancements is expected to yield:

\begin{itemize}[leftmargin=*]
    \item \textbf{Increased Accessibility:} 100\% keyboard navigability and screen reader compatibility
    \item \textbf{Improved Engagement:} 30-40\% increase in session duration through diverse interaction modalities
    \item \textbf{Enhanced Learning Outcomes:} 25\% improvement in symbol retention through adaptive feedback
    \item \textbf{Better Progress Monitoring:} Comprehensive session recordings enabling evidence-based instructional adjustments
    \item \textbf{Greater Parental Involvement:} Shareable session videos facilitating home-school collaboration
    \item \textbf{Reduced Sensory Overwhelm:} Customizable interface accommodating diverse sensory profiles
\end{itemize}

\subsection{Future Research Directions}

Several avenues for continued research and development emerge from this work:

\subsubsection{Machine Learning Integration}

\begin{itemize}[leftmargin=*]
    \item Predictive models for learning difficulty identification
    \item Automated adaptive difficulty adjustment
    \item Engagement prediction and intervention triggering
    \item Personalized learning path generation
\end{itemize}

\subsubsection{Advanced Assistive Technologies}

\begin{itemize}[leftmargin=*]
    \item Eye-tracking integration for gaze-based control
    \item Voice command recognition for hands-free operation
    \item Brain-computer interface exploration for severe motor impairments
    \item Haptic feedback systems for touch devices
\end{itemize}

\subsubsection{Extended Mathematical Content}

\begin{itemize}[leftmargin=*]
    \item Interactive problem-solving activities
    \item Step-by-step equation solving
    \item Real-world application scenarios
    \item Adaptive practice problem generation
\end{itemize}

\subsubsection{Social Learning Features}

\begin{itemize}[leftmargin=*]
    \item Collaborative problem-solving with peers
    \item Virtual learning companions
    \item Parent-child co-learning modes
    \item Teacher-facilitated group activities
\end{itemize}

\subsubsection{Longitudinal Studies}

\begin{itemize}[leftmargin=*]
    \item Long-term learning outcome assessment
    \item Skill generalization to real-world contexts
    \item Comparative studies with traditional instruction
    \item Cost-effectiveness analysis
\end{itemize}

\subsection{Limitations and Considerations}

Several limitations should be acknowledged:

\begin{enumerate}[leftmargin=*]
    \item \textbf{Individual Variability:} Autism spectrum diversity means no single approach works for all children; extensive customization remains necessary.
    
    \item \textbf{Technology Access:} Digital divide issues may limit access for underprivileged populations; offline capabilities and low-bandwidth optimization needed.
    
    \item \textbf{Screen Time Concerns:} Balancing benefits of technology-enhanced learning with risks of excessive screen exposure requires careful monitoring.
    
    \item \textbf{Human Interaction:} Technology supplements but cannot replace human instructors, therapists, and caregivers; blended approaches optimal.
    
    \item \textbf{Privacy Risks:} Screen recording features introduce data security and privacy concerns requiring robust safeguards.
\end{enumerate}

\subsection{Concluding Remarks}

The integration of evidence-based pedagogical strategies with advanced interaction technologies creates powerful opportunities for enhancing mathematical learning in children with Autism Spectrum Disorder. The proposed enhancements to Star Math Explorer—encompassing keyboard navigation, advanced mouse and touch events, engagement monitoring, and screen capture capabilities—promise to significantly improve accessibility, engagement, and learning outcomes.

By respecting the unique cognitive and sensory profiles of autism learners while leveraging modern web technologies, we can create inclusive educational experiences that not only accommodate differences but celebrate and capitalize upon the distinctive strengths children with ASD bring to mathematical thinking.

The future of autism education lies in personalized, adaptive, multimodal learning environments that provide structure without rigidity, challenge without overwhelm, and support without limitation. This work represents a step toward that vision, grounded in research evidence and empowered by technological innovation.

\vspace{1cm}

\begin{center}
\textit{"The only disability in life is a bad attitude."} \\
--- Scott Hamilton
\end{center}

\vspace{0.5cm}

\begin{center}
\textit{"Different, not less."} \\
--- Temple Grandin
\end{center}

\newpage

% ==================== REFERENCES ====================
\section*{References}
\addcontentsline{toc}{section}{References}

\begin{enumerate}[leftmargin=*]
    \item Baron-Cohen, S., Ashwin, E., Ashwin, C., Tavassoli, T., \& Chakrabarti, B. (2009). Talent in autism: Hyper-systemizing, hyper-attention to detail and sensory hypersensitivity. \textit{Philosophical Transactions of the Royal Society B: Biological Sciences, 364}(1522), 1377-1383.
    
    \item Burns, M. K., Codding, R. S., Boice, C. H., \& Lukito, G. (2010). Meta-analysis of acquisition and fluency math interventions with instructional and frustrational level skills: Evidence for a skill-by-treatment interaction. \textit{School Psychology Review, 39}(1), 69-83.
    
    \item Centers for Disease Control and Prevention. (2023). \textit{Data and Statistics on Autism Spectrum Disorder}. Retrieved from https://www.cdc.gov/autism/data/
    
    \item Goldstein, G., Johnson, C. R., \& Minshew, N. J. (2014). Attentional processes in autism. \textit{Journal of Autism and Developmental Disorders, 44}(10), 2474-2485.
    
    \item Grandin, T. (2009). \textit{Thinking in Pictures, Expanded Edition: My Life with Autism}. Vintage Books.
    
    \item Minshew, N. J., Meyer, J., \& Goldstein, G. (2002). Abstract reasoning in autism: A dissociation between concept formation and concept identification. \textit{Neuropsychology, 16}(3), 327-334.
    
    \item Pennington, R. C. (2010). Computer-assisted instruction for teaching academic skills to students with autism spectrum disorders: A review of literature. \textit{Focus on Autism and Other Developmental Disabilities, 25}(4), 239-248.
    
    \item React Documentation. (2026). \textit{React - The library for web and native user interfaces}. Retrieved from https://react.dev
    
    \item Schlosser, R. W., Shane, H. C., Allen, A. A., Abramson, J., Laubscher, E., \& Dimery, K. (2012). Just-in-time supports in augmentative and alternative communication. \textit{Seminars in Speech and Language, 33}(4), 298-322.
    
    \item Web Content Accessibility Guidelines (WCAG) 2.1. (2018). W3C Recommendation. Retrieved from https://www.w3.org/TR/WCAG21/
    
    \item Winter-Messiers, M. A., Herr, C. M., Wood, C. E., Brooks, A. P., Gates, M. A. M., Houston, T. L., \& Tingstad, K. I. (2007). How far can Brian ride the Daylight 4449 Express? A strength-based model of Asperger syndrome based on special interest areas. \textit{Focus on Autism and Other Developmental Disabilities, 22}(2), 67-79.
    
    \item Witzel, B. S., Mercer, C. D., \& Miller, M. D. (2003). Teaching algebra to students with learning difficulties: An investigation of an explicit instruction model. \textit{Learning Disabilities Research \& Practice, 18}(2), 121-131.
    
    \item Yakubova, G., Hughes, E. M., \& Shinaberry, M. (2015). Learning with technology: Video modeling with concrete-representational-abstract sequencing for students with autism spectrum disorder. \textit{Journal of Autism and Developmental Disorders, 46}(7), 2349-2362.
    
    \item MediaStream Recording API. (2026). MDN Web Docs. Retrieved from https://developer.mozilla.org/en-US/docs/Web/API/MediaStream\_Recording\_API
    
    \item React Screen Capture. (2026). NPM Package Documentation. Retrieved from https://www.npmjs.com/package/react-screen-capture
    
    \item COPPA - Children's Online Privacy Protection Act. (1998). Federal Trade Commission. Retrieved from https://www.ftc.gov/enforcement/rules/rulemaking-regulatory-reform-proceedings/childrens-online-privacy-protection-rule
    
    \item FERPA - Family Educational Rights and Privacy Act. (1974). U.S. Department of Education. Retrieved from https://www2.ed.gov/policy/gen/guid/fpco/ferpa/
    
    \item American Psychiatric Association. (2013). \textit{Diagnostic and Statistical Manual of Mental Disorders} (5th ed.). Arlington, VA: American Psychiatric Publishing.
\end{enumerate}

\newpage

% ==================== APPENDICES ====================
\appendix

\section{Code Repository Structure}

\begin{verbatim}
lab2/
├── frontend/
│   ├── src/
│   │   ├── components/
│   │   │   ├── Navigation.jsx
│   │   │   ├── StarField.jsx
│   │   │   ├── VoiceToggle.jsx
│   │   │   ├── ScreenRecorder.jsx (NEW)
│   │   │   ├── KeyboardNav.jsx (NEW)
│   │   │   └── SensoryControls.jsx (NEW)
│   │   ├── hooks/
│   │   │   ├── useKeyboardNavigation.js (NEW)
│   │   │   ├── useMouseTracking.js (NEW)
│   │   │   ├── useScreenCapture.js (NEW)
│   │   │   └── useAccessibility.js (NEW)
│   │   ├── pages/
│   │   │   ├── LandingPage.jsx
│   │   │   ├── LearnPage.jsx (ENHANCED)
│   │   │   └── ProgressPage.jsx (ENHANCED)
│   │   └── utils/
│   │       ├── voiceService.js
│   │       ├── eventTracking.js (NEW)
│   │       └── analytics.js (NEW)
│   └── package.json
├── backend/
│   ├── server.js (ENHANCED)
│   ├── routes/
│   │   ├── sessions.js (NEW)
│   │   └── analytics.js (NEW)
│   └── package.json
└── documentation/
    ├── autism_portal_documentation.tex
    ├── Star_Math_Explorer_Documentation.tex
    └── Enhanced_Interaction_Strategies_Report.tex (THIS DOCUMENT)
\end{verbatim}

\section{Installation and Setup Instructions}

\subsection{Prerequisites}

\begin{verbatim}
Node.js version 18.x or higher
npm or yarn package manager
Modern browser with ES6+ support
Git for version control
\end{verbatim}

\subsection{Installation Steps}

\begin{verbatim}
1. Clone repository:
   git clone https://github.com/9059Rohith/lab2.git
   cd lab2

2. Install frontend dependencies:
   cd frontend
   npm install
   npm install react-screen-capture
   npm install @react-aria/interactions

3. Install backend dependencies:
   cd ../backend
   npm install
   npm install multer express-fileupload

4. Start development servers:
   # Terminal 1 - Backend
   cd backend
   node server.js

   # Terminal 2 - Frontend
   cd frontend
   npm run dev

5. Access application:
   Navigate to http://localhost:5173
\end{verbatim}

\section{Accessibility Testing Checklist}

\begin{tcolorbox}[colback=lightgray, colframe=primaryblue, title=WCAG 2.1 Compliance Checklist]
\begin{enumerate}[leftmargin=*]
    \item[$\square$] All images have appropriate alt text
    \item[$\square$] Color contrast ratios meet minimum standards
    \item[$\square$] All functionality accessible via keyboard
    \item[$\square$] No keyboard traps present
    \item[$\square$] Focus indicators clearly visible
    \item[$\square$] Heading structure logical and sequential
    \item[$\square$] ARIA labels on interactive elements
    \item[$\square$] Forms have associated labels
    \item[$\square$] Error messages clearly communicated
    \item[$\square$] Time limits can be disabled/extended
    \item[$\square$] No flashing content >3 times per second
    \item[$\square$] Page title descriptive and unique
    \item[$\square$] Language of page specified
    \item[$\square$] Consistent navigation across pages
    \item[$\square$] Touch targets minimum 44x44 pixels
    \item[$\square$] Zoom to 200\% without horizontal scroll
    \item[$\square$] Content reflows at small viewport sizes
    \item[$\square$] Motion can be paused/disabled
    \item[$\square$] Screen reader compatible
    \item[$\square$] Text spacing can be adjusted
\end{enumerate}
\end{tcolorbox}

\section{Parent/Teacher User Guide}

\subsection{Getting Started}

\textbf{First-Time Setup:}
\begin{enumerate}[leftmargin=*]
    \item Create child profile with name
    \item Configure sensory preferences
    \item Review privacy settings and consent
    \item Test screen recording (optional)
    \item Set session duration preferences
\end{enumerate}

\subsection{Keyboard Shortcuts Reference}

\begin{table}[H]
\centering
\begin{tabular}{|l|l|}
\hline
\textbf{Key} & \textbf{Action} \\
\hline
Arrow Keys & Navigate between symbols \\
Enter / Space & Select focused symbol \\
1-9 & Quick select symbols 1-9 \\
Escape & Close modal / Return to main \\
Tab & Move to next interactive element \\
Shift + Tab & Move to previous element \\
R & Start/Stop recording \\
V & Toggle voice on/off \\
H & Show help overlay \\
\hline
\end{tabular}
\caption{Keyboard Shortcuts}
\end{table}

\subsection{Interpreting Progress Data}

\textbf{Key Metrics:}
\begin{itemize}[leftmargin=*]
    \item \textbf{Engagement Score:} Percentage of time actively interacting
    \item \textbf{Success Rate:} Proportion of correct symbol identifications
    \item \textbf{Attention Duration:} Average focus time before distraction
    \item \textbf{Symbols Mastered:} Number meeting mastery criterion
    \item \textbf{Struggle Points:} Symbols requiring additional support
\end{itemize}

\begin{figure}[H]
    \centering
    \includegraphics[width=0.75\textwidth]{figure-6.png}
    \caption{Proposed Enhancement - AR Symbol Recognition and Future Directions}
    \label{fig:future-enhancements}
\end{figure}

% ==================== END OF DOCUMENT ====================
\end{document}
